\documentclass[11pt,a4paper]{article}
\usepackage[margin=1in]{geometry}
\usepackage[utf8]{inputenc}
\usepackage[T1]{fontenc}
\usepackage{hyperref}
\usepackage{listings}
\usepackage{xcolor}
\usepackage{graphicx}
\usepackage{booktabs}
\usepackage{fancyhdr}
\usepackage{tabularx}
\usepackage{array}
\usepackage{amsmath}
\usepackage{colortbl}
\usepackage{setspace}

% ---- colours ----
\definecolor{codebg}{RGB}{245,245,245}
\definecolor{keywd}{RGB}{0,80,180}
\definecolor{commentcol}{RGB}{60,140,60}
\definecolor{stringcol}{RGB}{180,0,0}
\definecolor{notebg}{RGB}{255,248,220}
\definecolor{tipbg}{RGB}{230,245,255}
\definecolor{tabhead}{RGB}{220,230,242}

% ---- code listings ----
\lstset{
    basicstyle=\ttfamily\small,
    breaklines=true,
    frame=single,
    backgroundcolor=\color{codebg},
    keywordstyle=\color{keywd}\bfseries,
    commentstyle=\color{commentcol},
    stringstyle=\color{stringcol},
    upquote=true,
    numbers=none,
    xleftmargin=6pt,
    xrightmargin=6pt,
}

% ---- callout boxes ----
\newenvironment{notebox}%
    {\par\vspace{4pt}\noindent\begin{tabular}{|>{\columncolor{notebg}}p{0.96\linewidth}|}
     \hline \noindent\textbf{Note:} \\}%
    {\\ \hline \end{tabular}\par\vspace{4pt}}

\newenvironment{tipbox}%
    {\par\vspace{4pt}\noindent\begin{tabular}{|>{\columncolor{tipbg}}p{0.96\linewidth}|}
     \hline \noindent\textbf{Tip:} \\}%
    {\\ \hline \end{tabular}\par\vspace{4pt}}

% ---- header / footer ----
\pagestyle{fancy}
\fancyhf{}
\rhead{VASP Workflow --- User Manual}
\lhead{\leftmark}
\cfoot{\thepage}
\renewcommand{\headrulewidth}{0.4pt}

% ---- meta ----
\title{%
    \textbf{VASP Workflow}\\[4pt]
    \large User Manual\\[6pt]
    \normalsize Browser GUI $\cdot$ CLI Agents $\cdot$ SLURM HPC
}
\author{Keivan Esfarjani}
\date{June 2026 \quad --- \quad
      \url{https://github.com/KeivanS/VASP-Flow}}

% ----
\begin{document}
\maketitle
\tableofcontents
\newpage

% ----
\section{Introduction}
% ----

VASP Workflow is a framework for setting up, running, and analysing DFT
calculations with VASP. Two interfaces are provided:

\begin{itemize}
    \item \textbf{Browser GUI} (\texttt{vasp-gui.py}) --- a Flask web
          application that lets you configure calculations through a form,
          run each step from the browser, and view band-structure and DOS
          plots in real time.
    \item \textbf{CLI agents} (\texttt{vasp-agent.py} for workstations,
          \texttt{vasp-agent-slurm.py} for SLURM clusters) --- command-line
          tools that read a plain-text \texttt{instructions.txt} file and
          generate all INCAR, KPOINTS, POSCAR, POTCAR, and job scripts
          automatically.
\end{itemize}

Both interfaces share the same \texttt{modules/} back-end, so every feature
is available in either.

\subsection{Key features}

\begin{itemize}
    \item Functionals: PBEsol, PBE, R2SCAN, HSE06, RVV10, LDA, AM05
    \item Spin-orbit coupling (SOC) and collinear spin (ISPIN=2)
    \item Magnetic materials: correct MAGMOM per atom; spin-resolved band
          and DOS plots (spin-up / spin-down)
    \item GGA+U (Dudarev, LDAUTYPE=2) per element and orbital
    \item Automatic k-path detection from POSCAR lattice vectors
    \item Convergence testing (ENCUT and k-mesh sweeps)
    \item DOS with orbital projections via \texttt{sumo}
    \item Phonons via phonopy, Wannier90 interface, DFPT Born charges
    \item Smart KPAR/NCORE auto-defaults that scale with MPI task count
    \item SLURM workflow with dependency-chained job submission
\end{itemize}

\subsection{Requirements}

\begin{tabularx}{\linewidth}{@{}lX@{}}
\toprule
Package & Notes \\
\midrule
VASP 5.4+ & \texttt{vasp\_std}, \texttt{vasp\_ncl} (for SOC),
            optionally \texttt{vasp\_gam} \\
POTCAR library & PAW pseudopotential files in per-element folders \\
Python 3.8+ & Conda environment recommended \\
Flask & \texttt{pip install flask} --- GUI only \\
sumo $\geq$ 2.0 & \texttt{pip install sumo} --- band and DOS plots \\
MPI & \texttt{mpirun}, \texttt{mpiexec}, or \texttt{srun} (SLURM) \\
phonopy & optional; \texttt{conda install -c conda-forge phonopy} \\
wannier90.x & optional; \texttt{conda install -c conda-forge wannier90} \\
\bottomrule
\end{tabularx}

% ----
\section{Installation}
% ----

\begin{lstlisting}[language=bash]
# 1. Clone the repository
git clone https://github.com/KeivanS/VASP-Flow.git
cd VASP-Flow

# 2. Install Python dependencies
pip install flask sumo

# 3. Optional: phonons and Wannier support
conda install -c conda-forge phonopy wannier90

# 4. First-time configuration (creates site.env from template)
make setup
nano site.env       # set VASP paths, POTCAR directory, MPI command
\end{lstlisting}

\subsection{site.env --- platform configuration}

\texttt{site.env} is read by \texttt{make run} and exported to the GUI and
CLI agents.  Edit it once for your system:

\begin{lstlisting}
PYTHON          = python3
VASP_STD        = /path/to/vasp_std
VASP_NCL        = /path/to/vasp_ncl
VASP_GAM        = /path/to/vasp_gam       # optional
VASP_POTCAR_DIR = /path/to/potpaw_PBE
MPI_LAUNCH      = mpirun -np              # or: srun
MPI_NP          = 16
WANNIER90_X     = wannier90.x
\end{lstlisting}

% ----
\section{Browser GUI}
% ----

\subsection{Starting the GUI}

\begin{lstlisting}[language=bash]
make run          # launches on http://localhost:5001
\end{lstlisting}

\textbf{On a SLURM cluster:} run on the login node and open the GUI on
your laptop via SSH port-forwarding:

\begin{lstlisting}[language=bash]
# On your laptop: open the tunnel
ssh -L 5001:localhost:5001 username@cluster

# On the cluster login node
make run

# Open http://localhost:5001 in your laptop's browser
\end{lstlisting}

\subsection{Three-tab workflow}

\begin{enumerate}
    \item \textbf{Setup} --- configure the project, then click
          \emph{Generate Workflow}.
    \item \textbf{Workflow} --- run each calculation step; live terminal
          output is streamed to the browser.
    \item \textbf{Results} --- view band-structure and DOS plots; access
          OUTCAR analysis.
\end{enumerate}

To return to a previous project, use the \emph{Select project} dropdown
at the top of the page.  \textbf{$\to$ Workflow} jumps to the run tab;
\textbf{Edit \& Regenerate} re-opens Setup with all settings pre-filled.

\subsection{Setup tab --- field guide}

\subsubsection{Structure and POTCAR}

\begin{tabularx}{\linewidth}{@{}lX@{}}
\toprule
Field & Description \\
\midrule
Project name & Folder created in the working directory. \\
POSCAR & Paste the crystal structure, or click \emph{Load file}. \\
POTCAR directory & Path to your PAW library
    (e.g.\ \texttt{/opt/vasp/potpaw\_PBE.54}).
    If an element has multiple PAW variants (e.g.\ \texttt{Ga}
    vs \texttt{Ga\_d}), a chooser appears. \\
\bottomrule
\end{tabularx}

\subsubsection{Functional}

\begin{tabularx}{\linewidth}{@{}lX@{}}
\toprule
Choice & Description \\
\midrule
PBEsol & Optimised for solids; recommended starting point. \\
PBE & Standard GGA; widely used. \\
R2SCAN & Meta-GGA; good accuracy/cost balance. \\
HSE06 & Hybrid --- accurate band gaps, significantly slower. \\
RVV10 & van der Waals; for layered or molecular materials. \\
LDA / AM05 & Rarely needed. \\
\bottomrule
\end{tabularx}

\subsubsection{Spin / SOC}

\begin{tabularx}{\linewidth}{@{}lX@{}}
\toprule
Choice & Effect on INCAR \\
\midrule
None & Non-magnetic. \\
Collinear spin (ISPIN=2) & Sets \texttt{ISPIN = 2}; reveals the
    \emph{Magnetic moment / atom} field (see Section~\ref{sec:magnetic}). \\
SOC $\parallel$ z/x/y & Sets \texttt{LSORBIT = .TRUE.},
    \texttt{LNONCOLLINEAR = .TRUE.}, and \texttt{SAXIS};
    uses the \texttt{vasp\_ncl} binary. \\
\bottomrule
\end{tabularx}

\subsubsection{Other flags}

\begin{tabularx}{\linewidth}{@{}lX@{}}
\toprule
Field & Effect \\
\midrule
Hexagonal BZ & Correct $M$ point for hexagonal structures. \\
2D / slab & Sets $k_z = 1$ (Gamma-only out-of-plane). \\
GGA+U & Adds Hubbard $U$ rows; see Section~\ref{sec:ggau}. \\
Execution profile & Select \emph{SLURM HPC} to generate \texttt{sbatch}
    scripts; edit \texttt{profiles/slurm.json} first. \\
MPI tasks & Number of CPU cores per calculation. \\
\bottomrule
\end{tabularx}

\subsection{Task checkboxes}

\begin{tabularx}{\linewidth}{@{}llX@{}}
\toprule
Task & Directory & Notes \\
\midrule
Structure relaxation & \texttt{01\_relax} &
    IBRION=2; should run first if geometry is not pre-optimised. \\
SCF & \texttt{02\_scf} &
    Self-consistent; required before bands or DOS. \\
Band structure & \texttt{03\_bands} &
    Line-mode KPOINTS; see Section~\ref{sec:kpath}. \\
DOS & \texttt{04\_dos} &
    Dense k-mesh, LORBIT=11; orbital projections via sumo. \\
DFPT & \texttt{06\_dfpt} &
    Born charges and dielectric tensor (IBRION=8). \\
Phonons & \texttt{07\_phonons} &
    Phonopy finite-displacement; DFPT needed for NAC. \\
Wannier90 & \texttt{05\_wannier} &
    NSCF interface; see Section~\ref{sec:wannier}. \\
\bottomrule
\end{tabularx}

\begin{notebox}
Normal execution order: \textbf{relaxation $\to$ SCF $\to$ bands $\to$ DOS}.
Each step copies POSCAR/CHGCAR from the previous one automatically.
\end{notebox}

\subsection{Band structure options}

Leave the \emph{k-path} field blank to trigger \textbf{automatic detection}
from the POSCAR lattice vectors (see Section~\ref{sec:kpath}).
Enter a path manually to override, e.g.\ \texttt{G-M-K-G}.

\subsection{DOS options}

Add projection rows to request orbital-resolved DOS:

\begin{center}
Element: \texttt{Mo} \quad Orbitals: \textbf{d} \quad\quad
Element: \texttt{S}  \quad Orbitals: \textbf{p}
\end{center}

Sumo-dosplot produces a projected DOS PDF.
For spin-polarised calculations spin-up is plotted positive and spin-down
is mirrored below zero automatically.

\subsection{Workflow tab}

Each step shows three controls:

\begin{itemize}
    \item \textbf{Run} --- submit the step (\texttt{sbatch} on SLURM,
          \texttt{bash} otherwise).
    \item \textbf{Edit files} --- open INCAR, KPOINTS, or POSCAR in the
          browser.  A \texttt{.bak} backup is created on first edit.
    \item \textbf{Status badge} --- running / completed / warning.
\end{itemize}

\textbf{Run All Steps} chains every step sequentially and stops on failure.

\subsection{Results tab}

\begin{tabularx}{\linewidth}{@{}lX@{}}
\toprule
Panel & Contents \\
\midrule
Band structure &
    Spin-polarised: spin-up (blue) and spin-down (red dashed) on the same
    k-axis, Fermi level at zero.
    Non-magnetic: single-colour sumo plot. \\
Cumulative DOS &
    Stacked filled-area by element.
    Spin-polarised: up above zero, down mirrored below. \\
Projected DOS &
    Orbital-resolved via sumo; spin-down mirrored automatically. \\
OUTCAR analysis &
    Energy, Fermi level, total magnetic moment, forces. \\
\bottomrule
\end{tabularx}

% ----
\section{CLI Agents}
% ----

\subsection{vasp-agent.py --- workstation}

\begin{lstlisting}[language=bash]
# Run from the directory containing instructions.txt and POSCAR
./vasp-agent.py
./vasp-agent.py -i my_inst.txt -s struct.vasp

# After generation
cd ProjectName
./run_all.sh          # runs every step in sequence
./analyze.sh          # plots after all steps complete
\end{lstlisting}

\subsection{vasp-agent-slurm.py --- SLURM clusters}

\begin{lstlisting}[language=bash]
./vasp-agent-slurm.py                  # uses profiles/slurm.json
./vasp-agent-slurm.py --profile slurm

cd ProjectName
./submit_all.sh        # sbatch with dependency chaining
squeue -u $USER        # monitor jobs
./analyze.sh           # run locally after jobs finish
\end{lstlisting}

Two-phase workflow when convergence tests are requested:

\begin{lstlisting}[language=bash]
./submit_convergence.sh     # STEP 1
# review results, then:
./submit_calculations.sh    # STEP 2 -- prompts for ENCUT/k-mesh
\end{lstlisting}

\subsection{Generated directory structure}

\begin{lstlisting}
ProjectName/
---- instructions.txt
---- POSCAR
---- POTCAR               built once from $VASP_POTCAR_DIR
---- 01_relax/            INCAR  KPOINTS  POSCAR  POTCAR  run.sh
---- 02_scf/
---- 03_bands/
---- 04_dos/
---- analysis/            plots generated by analyze.sh
---- run_all.sh           workstation sequential runner
---- submit_all.sh        SLURM job submission (dependency chained)
---- analyze.sh           post-processing: plots + OUTCAR extraction
\end{lstlisting}

% ----
\section{instructions.txt Format}
\label{sec:instructions}
% ----

\subsection{Basic structure}

\begin{lstlisting}
Project: My material name

Methods: PBEsol functional
         collinear magnetization
         GGA+U with U=4.0 on Fe-d orbitals

Tasks: structure relaxation
       SCF calculation
       band structure
       DOS (projected on Fe-d and O-p)

Convergence: Test k-points 6x6x6 to 12x12x12
             ENCUT 400-600 eV

MAGMOM: 4.0
NSW: 100
EDIFFG: -0.01
MPI: 128
NODES: 2
NTASKS_PER_NODE: 64
WALLTIME: 12:00:00
\end{lstlisting}

The parser is case-insensitive and searches the entire file, so section
order is flexible.

\subsection{Keyword reference}

\begin{tabularx}{\linewidth}{@{}lX@{}}
\toprule
Keyword & Meaning / example \\
\midrule
\texttt{Project:} & Project name; becomes the output folder name. \\
\texttt{MAGMOM:} &
    Initial moment in $\mu_{\mathrm{B}}$ per atom (collinear only).
    Written as \texttt{MAGMOM = N*value}. Default 2.0. \\
\texttt{ISIF:} &
    Relaxation type: 3=full, 2=ions only, 4=fixed volume,
    7=fixed shape, 6=volume only.  Default 3. \\
\texttt{NSW:} & Maximum ionic steps. Default 100. \\
\texttt{EDIFFG:} & Force criterion in eV/\AA{}. Default $-0.01$. \\
\texttt{ENCUT:} & Manual plane-wave cutoff override in eV. \\
\texttt{NKPTS:} & k-points per band segment. Default 40. \\
\texttt{MPI:} & Total MPI tasks (nodes $\times$ cores/node). \\
\texttt{KPAR:} & Global k-point parallelisation override. \\
\texttt{NCORE:} & Global cores-per-orbital override. \\
\texttt{SCF\_KPAR:} etc. &
    Per-step KPAR: also BANDS\_, DOS\_, RELAX\_, DFPT\_, PHONONS\_. \\
\texttt{SCF\_NCORE:} etc. &
    Per-step NCORE; same prefixes as above. \\
\texttt{NODES:} & Compute nodes (SLURM). \\
\texttt{NTASKS\_PER\_NODE:} & MPI tasks per node (SLURM). \\
\texttt{PARTITION:} & SLURM partition name. \\
\texttt{WALLTIME:} & Job time limit, e.g.\ \texttt{12:00:00}. \\
\texttt{ACCOUNT:} & SLURM account / allocation. \\
\texttt{DFPT\_EDIFF:} & EDIFF for DFPT step. Default \texttt{1E-8}. \\
\texttt{PHONONS\_DIM:} & Supercell dimensions, e.g.\ \texttt{2 2 2}. \\
\texttt{PHONONS\_MESH:} & q-mesh, e.g.\ \texttt{20 20 20}. \\
\texttt{PHONONS\_DISP:} & Displacement in \AA{}. Default 0.01. \\
\texttt{PHONONS\_BAND:} & High-symmetry path for phonon bands. \\
\texttt{PHONONS\_NAC:} & \texttt{FALSE} to disable NAC. Default TRUE. \\
\texttt{WANNIER\_NUM\_WANN:} & Number of Wannier functions. \\
\texttt{WANNIER\_PROJ:} & Projections, e.g.\ \texttt{Ga:s;As:p}. \\
\texttt{WANNIER\_EWIN:} & Energy window, e.g.\ \texttt{-5.0 to 15.0}. \\
\bottomrule
\end{tabularx}

% ----
\section{Methods}
% ----

\subsection{Functionals}

Write in the \texttt{Methods:} block:

\begin{lstlisting}
Methods: PBEsol functional
Methods: R2SCAN functional
Methods: HSE06 functional
Methods: RVV10 functional     # van der Waals
\end{lstlisting}

Default ENCUT values: PBE 400 eV, PBEsol 450 eV, R2SCAN 680 eV,
HSE06 400 eV, RVV10 450 eV.

\subsection{Spin-orbit coupling}

\begin{lstlisting}
Methods: SOC with magnetization in z-direction
         SOC with magnetization in x-direction
\end{lstlisting}

Uses \texttt{vasp\_ncl}.  Sets \texttt{LSORBIT = .TRUE.},
\texttt{LNONCOLLINEAR = .TRUE.}, and the appropriate \texttt{SAXIS} vector.

\subsection{Magnetic materials (collinear spin)}
\label{sec:magnetic}

\begin{lstlisting}
Methods: collinear magnetization
MAGMOM: 4.0
\end{lstlisting}

The generator sets \texttt{ISPIN = 2} and
\texttt{MAGMOM = $N \times \text{value}$} in every INCAR, where $N$ is
the total atom count read from the POSCAR.

\begin{tipbox}
Typical starting moments: Fe, Co $\approx$ 4--5 $\mu_{\mathrm{B}}$;
Ni $\approx$ 2 $\mu_{\mathrm{B}}$;
Cr, Mn $\approx$ 3--4 $\mu_{\mathrm{B}}$;
rare earths $\approx$ 7 $\mu_{\mathrm{B}}$.
The value only needs to be in the right ball-park; VASP converges
the self-consistent moment.
\end{tipbox}

\textbf{Antiferromagnets:} set a uniform \texttt{MAGMOM} in instructions,
then manually edit the INCAR to alternate signs:
\begin{lstlisting}
MAGMOM = 4.0 -4.0 4.0 -4.0   ! alternating per sublattice
\end{lstlisting}

\textbf{Spin-resolved plots produced automatically:}

\begin{itemize}
    \item \emph{Cumulative DOS} --- spin-up above the x-axis (positive);
          spin-down mirrored below zero.
          Elements are colour-coded; solid lines = up, dashed = down.
    \item \emph{Band structure} --- spin-up in blue, spin-down in red
          dashed, both on the same k-axis with Fermi level at zero.
    \item \emph{Projected DOS} --- sumo natively negates the spin-down
          channel, so the orbital-resolved PDF also shows a mirrored layout.
\end{itemize}

Check the final moment after calculation:
\begin{lstlisting}[language=bash]
grep "magnetization (x)" OUTCAR | tail -5
\end{lstlisting}

\subsection{GGA+U}
\label{sec:ggau}

\begin{lstlisting}
Methods: GGA+U with U=4.0 on Fe-d orbitals
         GGA+U with U=3.0 on Ni-d orbitals
         GGA+U with U=6.0 on Ce-f orbitals
\end{lstlisting}

Uses Dudarev scheme (LDAUTYPE=2); $J = 0$ is set automatically.
Can be combined with \texttt{collinear magnetization}.

In the GUI: check \emph{GGA+U} and add one row per element/orbital.

% ----
\section{K-Path for Band Structure}
\label{sec:kpath}
% ----

\subsection{Auto-detection (recommended)}

Omit the path in \texttt{instructions.txt}, or leave the field blank
in the GUI:

\begin{lstlisting}
Tasks: band structure        # no "along ..." -- auto-detect
\end{lstlisting}

The generator reads POSCAR lattice vectors, computes cell lengths and
angles, classifies the Bravais lattice, and selects the standard path.

\begin{tabularx}{\linewidth}{@{}llllX@{}}
\toprule
Type & Name & Detection criterion & Default path \\
\midrule
cF & FCC & $a{=}b{=}c$, angles $\approx 60^{\circ}$
        & G-X-W-K-G-L-U-W \\
cI & BCC & $a{=}b{=}c$, angles $\approx 109.5^{\circ}$
        & G-H-N-G-P-H \\
cP & Cubic & $a{=}b{=}c$, $90^{\circ}$
        & G-X-M-G-R-X \\
hP & Hexagonal & $a{\approx}b$, $\gamma{=}120^{\circ}$
        & G-M-K-G-A-L-H-A \\
tP & Tetragonal & $a{\approx}b{\neq}c$, $90^{\circ}$
        & G-X-M-G-Z-R-A-Z \\
oP & Orthorhombic & $90^{\circ}$
        & G-X-S-Y-G-Z-U-R-T-Z \\
rP & Rhombohedral & $a{=}b{=}c$, equal angles
        & G-T-F-G-L \\
mP & Monoclinic & $\alpha{=}\gamma{=}90^{\circ}$, $\beta{\neq}90^{\circ}$
        & G-Y-A-Z-G-B-D \\
\bottomrule
\end{tabularx}

\begin{notebox}
For \textbf{conventional} FCC or BCC cells (right angles, 4 or 2 atoms),
auto-detection returns the cubic (cP) path.
Use the \textbf{primitive} cell for an unfolded FCC/BCC band structure ---
the primitive FCC cell has all angles $\approx 60^{\circ}$ and is
recognised automatically.
\end{notebox}

The detected type is written into the KPOINTS comment:
\begin{lstlisting}
Line-mode KPOINTS  [FCC -- auto-detected]
\end{lstlisting}

\subsection{Manual override}

\begin{lstlisting}
Tasks: band structure along G-X-W-K-G-L
\end{lstlisting}

Available labels and fractional reciprocal coordinates:

\begin{tabularx}{\linewidth}{@{}llX@{}}
\toprule
Label & Coords & Notes \\
\midrule
G & (0, 0, 0) & $\Gamma$ \\
X & (0.5, 0, 0) & generic / hex M in non-hex \\
M & (0.5, 0.5, 0) & cubic, tetragonal \\
K & (1/3, 1/3, 0) & hexagonal \\
R & (0.5, 0.5, 0.5) & cubic \\
Z & (0, 0, 0.5) & general \\
A & (0, 0, 0.5) & hex / tet \\
L & (0.5, 0, 0.5) & hexagonal \\
H & (1/3, 1/3, 0.5) & hexagonal \\
S & (0.5, 0.5, 0) & orthorhombic \\
T & (0, 0.5, 0.5) & orthorhombic \\
U & (0.5, 0, 0.5) & orthorhombic \\
Y & (0, 0.5, 0) & ortho / monoclinic \\
W & (0.5, 0.25, 0.75) & \textbf{FCC primitive} \\
P & (0.25, 0.25, 0.25) & \textbf{BCC primitive} \\
N & (0, 0, 0.5) & \textbf{BCC primitive} \\
\bottomrule
\end{tabularx}

% ----
\section{DOS Projections}
% ----

All of the following forms are accepted:

\begin{lstlisting}
DOS (projected on Mo-d and S-p)
DOS with projections on Ni:d and Ni:s
DOS with projection on Fe-d and O-p
\end{lstlisting}

Element--orbital syntax: \texttt{El-orbital} or \texttt{El:orbital};
orbital channels are \texttt{s}, \texttt{p}, \texttt{d}, \texttt{f}.

Sumo-dosplot generates an orbital-resolved PDF in \texttt{analysis/}.
For spin-polarised calculations, spin-up is positive and spin-down is
negative (automatic).

% ----
\section{Convergence Tests}
% ----

\begin{lstlisting}
# Range syntax
Convergence: Test k-points from 6x6x1 to 18x18x1    # 2D
             Test k-points from 6x6x6 to 12x12x12   # 3D
             ENCUT 400-600 eV

# Explicit mesh list
Convergence: Test k-points 6x6x3, 12x12x6, 18x18x9
\end{lstlisting}

Results in \texttt{00\_convergence/kpoints/} and
\texttt{00\_convergence/encut/}.

\begin{lstlisting}[language=bash]
./run_convergence.sh          # workstation
./submit_convergence.sh       # SLURM

./analyze_convergence.sh      # plots energy / pressure / forces
\end{lstlisting}

% ----
\section{Parallelization}
% ----

\subsection{Auto-defaults}

When KPAR and NCORE are not specified, the agent computes them from
$N$ = total MPI tasks:

\begin{tabularx}{\linewidth}{@{}llll@{}}
\toprule
Step & KPAR & NCORE & Constraint \\
\midrule
SCF / DOS / relax & $\lfloor\sqrt{N}\rfloor$ & $N/\mathrm{KPAR}$ & \\
bands & 1 & $N$ & KPAR=1 for line-mode KPOINTS \\
DFPT & 1 & $N$ & KPAR=1 required for IBRION=8 \\
phonons & 1 & $\lfloor\sqrt{N}\rfloor$ & \\
\bottomrule
\end{tabularx}

KPAR $\times$ NCORE must divide $N$ evenly; the agent rounds down automatically.

\subsection{Manual overrides}

\begin{lstlisting}
MPI: 128             # total tasks
KPAR: 4              # global
NCORE: 8             # global

SCF_KPAR: 4          # per-step (take priority)
BANDS_KPAR: 1
DOS_KPAR: 8
RELAX_KPAR: 2
DFPT_KPAR: 1
PHONONS_KPAR: 1

SCF_NCORE: 8
BANDS_NCORE: 16
DOS_NCORE: 8
\end{lstlisting}

% ----
\section{SLURM Cluster Settings}
% ----

\subsection{profiles/slurm.json}

\begin{lstlisting}
{
  "vasp_std": "vasp_std",
  "vasp_ncl": "vasp_ncl",
  "mpi_cmd":  "srun",
  "modules":  ["intel/2021", "impi/2021", "vasp/6.3.2"],
  "slurm": {
    "partition":       "standard",
    "nodes":           2,
    "ntasks_per_node": 64,
    "time":            "12:00:00",
    "account":         "mygroup",
    "output":          "slurm-%j.out",
    "error":           "slurm-%j.err"
  }
}
\end{lstlisting}

\subsection{Per-project overrides in instructions.txt}

\begin{lstlisting}
NODES: 4
NTASKS_PER_NODE: 32
PARTITION: gpu
WALLTIME: 48:00:00
ACCOUNT: myallocation
\end{lstlisting}

% ----
\section{Wannier90}
\label{sec:wannier}
% ----

\begin{lstlisting}
Tasks: Wannier90 wannierization

WANNIER_NUM_WANN: 8
WANNIER_PROJ: Ga:s;As:p
WANNIER_EWIN: -5.0 to 15.0
\end{lstlisting}

The agent generates \texttt{05\_wannier/wannier90.win} with the NSCF k-mesh
and projection block.  After the NSCF run, execute \texttt{wannier90.x}
from the \texttt{05\_wannier/} directory.

% ----
\section{Phonons (phonopy)}
% ----

\begin{lstlisting}
Tasks: phonons lattice dynamics

PHONONS_DIM: 2 2 2
PHONONS_MESH: 20 20 20
PHONONS_DISP: 0.01
PHONONS_BAND: G M K G A
PHONONS_NAC: FALSE          # disable non-analytic correction
\end{lstlisting}

\begin{notebox}
For the non-analytic correction (important for polar materials), include
the \texttt{DFPT} task in your list so Born effective charges are available.
If phonopy is the only post-SCF step, set \texttt{PHONONS\_NAC: FALSE}.
\end{notebox}

% ----
\section{Worked Examples}
% ----

\subsection{MoS\textsubscript{2} monolayer with SOC}

\begin{lstlisting}
Project: MoS2 monolayer

Methods: PBEsol functional
         SOC with magnetization in z-direction

Tasks: structure relaxation
       SCF calculation
       band structure
       DOS (projected on Mo-d and S-p)

Convergence: Test k-points 8x8 to 20x20
             ENCUT 400-600 eV

MPI: 64
\end{lstlisting}

Omitting the k-path lets the generator detect hexagonal (hP) from the POSCAR
and use G-M-K-G-A-L-H-A automatically.

\subsection{Nickel: collinear magnetism with GGA+U}

\begin{lstlisting}
Project: Nickel GGA+U

Methods: PBEsol functional
         collinear magnetization
         GGA+U with U=4.0 on Ni-d orbitals

Tasks: structure relaxation
       SCF calculation
       band structure
       DOS (projected on Ni-d and Ni-s)

MAGMOM: 2.0
MPI: 128
NODES: 2
NTASKS_PER_NODE: 64
WALLTIME: 08:00:00
\end{lstlisting}

With the primitive FCC POSCAR (cell angles $\approx 60^{\circ}$), the
generator detects cF and uses G-X-W-K-G-L-U-W.
Band and DOS plots show both spin channels separately.

\subsection{GaAs: phonons with NAC}

\begin{lstlisting}
Project: GaAs phonons

Methods: PBEsol functional

Tasks: SCF calculation
       DFPT Born charges and dielectric
       phonons lattice dynamics

PHONONS_DIM: 4 4 4
PHONONS_MESH: 20 20 20
PHONONS_BAND: G X W L G K

MPI: 64
WALLTIME: 24:00:00
\end{lstlisting}

DFPT is included so Born effective charges are available for the NAC.
The phonon band path is specified manually.

\subsection{Bi\textsubscript{2}Se\textsubscript{3} topological insulator}

\begin{lstlisting}
Project: Bi2Se3 topological

Methods: PBEsol functional
         SOC with magnetization in z-direction

Tasks: structure relaxation
       SCF calculation
       band structure
       DOS (projected on Bi-p and Se-p)
       Wannier90 wannierization

Convergence: Test k-points from 8x8x4 to 16x16x8
             ENCUT 400-600 eV

WANNIER_NUM_WANN: 18
WANNIER_PROJ: Bi:p;Se:p
WANNIER_EWIN: -2.0 to 2.0
MPI: 128
\end{lstlisting}

% ----
\section{Common OUTCAR Checks}
% ----

\begin{lstlisting}[language=bash]
# Converged?
grep "reached required accuracy"  OUTCAR

# Total energy (last ionic step)
grep "energy  without"            OUTCAR | tail -1

# Fermi level
grep "E-fermi"                    OUTCAR | tail -1

# Magnetic moment (collinear)
grep "magnetization (x)"          OUTCAR | tail -5

# Forces
grep "TOTAL-FORCE"                OUTCAR | tail -20

# Born effective charges (after DFPT)
grep -A3 "BORN EFFECTIVE CHARGES" OUTCAR | head -20
\end{lstlisting}

% ----
\section{Troubleshooting}
% ----

\begin{tabularx}{\linewidth}{@{}lX@{}}
\toprule
Symptom & Solution \\
\midrule
POTCAR not found &
    Check \texttt{VASP\_POTCAR\_DIR} in \texttt{site.env} or the GUI
    System Setup.  For unusual PAW variants (\texttt{Ga\_d},
    \texttt{Fe\_sv}), select the correct one in the POTCAR chooser. \\[4pt]
\texttt{vasp\_std} not found &
    Set the full path in \texttt{site.env} or GUI System Setup. \\[4pt]
WARNING: may not have converged &
    Open OUTCAR; increase NSW (ionic steps) or NELM (electronic steps)
    in INCAR, then re-run. \\[4pt]
Wrong band structure ($W = \Gamma$) &
    The POSCAR is a conventional cubic cell.  Use the primitive cell for
    FCC/BCC so auto-detection picks cF/cI. \\[4pt]
Projected DOS missing &
    Check instructions.txt uses \emph{projected on} or
    \emph{projections on}; both are accepted.
    Verify element names match POSCAR. \\[4pt]
KPAR error &
    KPAR $\times$ NCORE must divide total MPI tasks.
    Let the agent use auto-defaults, or set valid values explicitly. \\[4pt]
GUI not reachable on laptop &
    Check the SSH tunnel:
    \texttt{ssh -L 5001:localhost:5001 user@cluster}. \\[4pt]
Project missing from dropdown &
    Folder must contain both \texttt{POSCAR} and \texttt{settings.json}
    and be in the directory where \texttt{vasp-gui.py} was started. \\
\bottomrule
\end{tabularx}

% ----
\section{Version History}
% ----

\begin{description}
    \item[v1.0 (2026-03-23)] Initial release --- GUI, workstation agent,
        PBEsol/R2SCAN/HSE06, SOC, GGA+U, convergence tests,
        Wannier90, phonopy.
    \item[v1.1 (2026-04--05)] Per-task KPAR/NCORE smart defaults;
        SLURM agent with dependency-chained submission;
        cumulative DOS by element (no pymatgen);
        copy-from-relax/scf scripts.
    \item[v1.2 (2026-06)] Collinear magnetism: MAGMOM per atom
        from POSCAR, spin-resolved band and DOS plots.
        Auto k-path detection from POSCAR Bravais lattice
        (8 crystal systems).
        Fixed DOS projection parsing.
        Separate GUI and agent quick-reference sheets.
\end{description}

\vfill\hrule\vspace{6pt}
\noindent\small
Source: \url{https://github.com/KeivanS/VASP-Flow}\hfill
Compiled: \today\hfill
Maintained by Keivan Esfarjani

\end{document}
